% Mathématiques 3e — V3 LaTeX
% Calcul littéral, équations et inéquations

\section{Objectifs}
\begin{itemize}
\item Développer et factoriser des expressions avec la double distributivité.
\item Distinguer identité et équation et démontrer des propriétés littérales.
\item Résoudre des équations du premier degré et des équations produits simples.
\item Résoudre et interpréter des inéquations dans des problèmes.
\end{itemize}

\section{Repères et vocabulaire}
Le calcul littéral devient en troisième un langage de modélisation et de preuve. Développer ou factoriser n’est pas une fin en soi : on transforme une expression parce qu’une forme est plus utile qu’une autre pour calculer, démontrer, résoudre une équation ou comparer deux situations.

\section{1. Structure d’une expression}
Avant de calculer, on identifie si l’expression est une somme ou un produit. Cette lecture évite de distribuer ou de factoriser au mauvais endroit.

\[
3(x+5)$ est un produit
\]

\[
3x+15$ est une somme
\]

\begin{exemple}
Deux expressions peuvent être égales pour toute valeur de $x$ tout en ayant des structures différentes.
\end{exemple}

\section{2. Double distributivité}
La double distributivité découle de deux applications successives de la distributivité simple. Elle permet de développer un produit de deux sommes.

\[
(a+b)(c+d)=ac+ad+bc+bd
\]

\begin{exemple}
On peut démontrer cette identité en posant temporairement $k=a+b$ puis en utilisant deux fois la distributivité.
\end{exemple}

\coursefigure{/images/mathematiques/3eme/calcul-litteral-equations-inequations-1.svg}{Représentation structurante pour le chapitre Calcul littéral, équations et inéquations.}{Cette figure organise visuellement les relations utilisées dans les premières notions du chapitre.}

\section{3. Factoriser}
Factoriser consiste à transformer une somme en produit. On recherche un facteur commun ou une identité adaptée lorsque le programme le permet.

\[
6x+18=6(x+3)
\]

\[
a^2-b^2=(a-b)(a+b)
\]

\begin{exemple}
La factorisation est utile pour résoudre une équation produit ou faire apparaître une propriété générale.
\end{exemple}

\section{4. Identité et équation}
Une identité est vraie pour toutes les valeurs admissibles ; une équation n’est vraie que pour certaines valeurs appelées solutions.

\[
(x+1)^2-x^2=2x+1
\]

\begin{exemple}
Tester quelques valeurs peut aider à conjecturer une identité, mais seule une transformation générale la démontre.
\end{exemple}

\section{5. Équation du premier degré}
Résoudre une équation consiste à conserver l’égalité tout en isolant l’inconnue. Chaque opération effectuée sur un membre doit être effectuée sur l’autre.

\[
7x+4=3x+28\Rightarrow x=6
\]

\begin{exemple}
La vérification remplace la solution dans l’équation initiale et permet de détecter une erreur de signe ou de calcul.
\end{exemple}

\coursefigure{/images/mathematiques/3eme/calcul-litteral-equations-inequations-2.svg}{Deuxième représentation pédagogique pour le chapitre Calcul littéral, équations et inéquations.}{Cette représentation sert de support de lecture, de comparaison ou de contrôle dans les problèmes du chapitre.}

\section{6. Équation produit}
Un produit est nul si et seulement si au moins un de ses facteurs est nul. La factorisation transforme alors certaines équations en plusieurs équations simples.

\[
(x-4)(x+2)=0\Rightarrow x=4\text{ ou }x=-2
\]

\begin{exemple}
Pour une équation comme $x^2=49$, on peut écrire $x^2-7^2=0$ puis factoriser.
\end{exemple}

\section{7. Inéquation}
Une inéquation cherche un ensemble de valeurs vérifiant une comparaison. On représente souvent les solutions sur une droite graduée.

\[
3x+5<20\Rightarrow x<5
\]

\begin{exemple}
Lorsque l’on multiplie ou divise une inégalité par un nombre négatif, le sens de l’inégalité change.
\end{exemple}

\section{8. Modéliser un choix}
Comparer deux tarifs conduit souvent à une équation pour le seuil d’égalité puis à une inéquation pour déterminer quelle offre est avantageuse.

\[
12+4x<5{,}5x\Rightarrow x>8
\]

\begin{exemple}
Une solution algébrique doit être interprétée : si $x$ compte des séances, on ne conserve que les valeurs entières compatibles avec la situation.
\end{exemple}

\section{Méthode générale de résolution}
\begin{methode}
\begin{enumerate}
\item Nommer clairement l’inconnue.
\item Traduire la situation par une expression, une équation ou une inéquation.
\item Réduire ou factoriser seulement si cela simplifie le problème.
\item Effectuer des transformations équivalentes sur les deux membres.
\item Vérifier les solutions dans l’égalité ou la contrainte initiale.
\item Interpréter les valeurs admissibles dans le contexte.
\end{enumerate}
\end{methode}

\section{Problème guidé et stratégie}
Dans un problème de troisième, la difficulté ne vient pas seulement du calcul. Il faut reconnaître la notion utile, choisir une représentation, enchaîner plusieurs étapes et expliquer pourquoi la méthode est valide. On commence donc par isoler les données, puis on écrit la relation mathématique avant de remplacer par des nombres. Une réponse est considérée comme complète lorsqu’elle comporte une justification, un calcul lisible, un contrôle et une interprétation.

\begin{attention}
Une figure, un graphique, un tableau ou une calculatrice fournit des informations, mais ne remplace pas une justification lorsque le problème demande une preuve. Inversement, un calcul exact sans interprétation peut rester incomplet dans un contexte concret.
\end{attention}

\section{Contrôler un résultat}
Avant de valider une réponse, on vérifie le signe, l’ordre de grandeur, les unités et la compatibilité avec les hypothèses. On peut aussi refaire le calcul par une autre représentation : formule et graphique, valeur exacte et approximation, calcul direct et vérification dans l’énoncé. Cette habitude permet de repérer une erreur de saisie ou une mauvaise modélisation avant la conclusion.

\section{Erreurs fréquentes}
\begin{itemize}
\item Développer un seul terme d’une parenthèse.
\item Confondre identité et équation.
\item Changer un signe sans opération équivalente.
\item Oublier les deux solutions possibles de $x^2=a$ lorsque $a>0$.
\item Oublier d’inverser le sens lors d’une multiplication ou division par un nombre négatif.
\item Donner une solution numérique incompatible avec le contexte.
\end{itemize}

\section{À retenir}
\begin{aretenir}
\begin{itemize}
\item La forme développée et la forme factorisée répondent à des besoins différents.
\item Une identité est générale ; une équation sélectionne des valeurs.
\item Un produit nul conduit à l’annulation d’au moins un facteur.
\item Une inéquation décrit un ensemble de valeurs et doit être interprétée.
\end{itemize}
\end{aretenir}

\section{Réinvestissement : Équation produit}
Pour réinvestir cette notion, on part d’une situation où plusieurs représentations sont possibles. Un produit est nul si et seulement si au moins un de ses facteurs est nul. La factorisation transforme alors certaines équations en plusieurs équations simples. L’élève doit expliciter son choix, effectuer le calcul et revenir au contexte. Pour une équation comme $x^2=49$, on peut écrire $x^2-7^2=0$ puis factoriser. Le contrôle final consiste à vérifier que la réponse respecte les hypothèses et que son ordre de grandeur est compatible avec les données.

\section{Réinvestissement : Structure d’une expression}
Pour réinvestir cette notion, on part d’une situation où plusieurs représentations sont possibles. Avant de calculer, on identifie si l’expression est une somme ou un produit. Cette lecture évite de distribuer ou de factoriser au mauvais endroit. L’élève doit expliciter son choix, effectuer le calcul et revenir au contexte. Deux expressions peuvent être égales pour toute valeur de $x$ tout en ayant des structures différentes. Le contrôle final consiste à vérifier que la réponse respecte les hypothèses et que son ordre de grandeur est compatible avec les données.
